% Shared preamble of the TikZ figures (tools/figures.py). Labels are mathematics only, because the
% figures serve both languages. Words go into the captions.
% Placeholder colours are mapped to the page colours in book/assets/style.css (light and dark theme):
% black = text colour, acc = accent (blue), nov = state after loading (orange), siv = auxiliary (grey).
\usetikzlibrary{arrows.meta, calc, decorations.pathreplacing}
\definecolor{acc}{HTML}{0000FF}
\definecolor{nov}{HTML}{FF0000}
\definecolor{siv}{HTML}{00FF00}
\providecommand{\vek}[1]{\mathbf{#1}}
\tikzset{
  >={Stealth[length=2.4mm, width=1.7mm]},
  every picture/.style={line width=0.6pt, line cap=round, line join=round},
  dim/.style={siv, <->, line width=0.45pt},
  aux/.style={siv, dashed, line width=0.45pt},
  brace/.style={siv, line width=0.45pt, decorate, decoration={brace, amplitude=4pt}},
  dot/.style={circle, fill, inner sep=1.3pt},
  hole/.style={circle, draw, fill=white, inner sep=1.6pt},
  hatch/.style={siv, line width=0.45pt},
}
% hatching above a horizontal line from (#1) of length #2: \hatchup{x,y}{len}
\newcommand{\hatchup}[2]{\draw[hatch] (#1) -- ++(#2,0);
  \foreach \i in {0,0.25,...,#2} {\draw[hatch] ($(#1)+(\i,0)$) -- ++(0.2,0.2);}}
\newcommand{\hatchdown}[2]{\draw[hatch] (#1) -- ++(#2,0);
  \foreach \i in {0,0.25,...,#2} {\draw[hatch] ($(#1)+(\i,0)$) -- ++(-0.2,-0.2);}}
\newcommand{\hatchleft}[2]{\draw[hatch] (#1) -- ++(0,#2);
  \foreach \i in {0,0.25,...,#2} {\draw[hatch] ($(#1)+(0,\i)$) -- ++(-0.2,-0.2);}}
